\documentclass[10pt,a4paper]{amsart}
\usepackage[margin=19mm]{geometry}
\usepackage{amsmath,amssymb,mathtools}
\usepackage[scr=rsfs,cal=euler]{mathalfa}
\usepackage{xfrac}
\usepackage[unicode,hidelinks,bookmarks=false]{hyperref}
\newcommand{\ie}{i.e.\ }
\newcommand{\cf}{cf.\ }
\newcommand{\resp}{respectively\ }
\newcommand{\Subsection}{\subsection}
\providecommand{\nobreakdash}{\nobreak\hbox{-}}
\setlength{\emergencystretch}{2em}
\multlinegap0pt
\usepackage{amssymb}
\usepackage{mathtools}
\usepackage{float}
\usepackage[shortlabels]{enumitem}
\usepackage{csquotes}
\newtheorem{definition}{Definition}
\newtheorem{proposition}{Proposition}
\newtheorem{theorem}{Theorem}
\newcommand{\op}{\operatorname}
\title{Trapezoidality of flat arrangement polynomials}
\author{Yuan Gao, Tianyu Yuan}
\date{Source: arXiv:2608.21855v1 (CC BY 4.0)}
\begin{document}
\maketitle

\noindent\emph{Source and licence.} Trapezoidality of flat arrangement polynomials. Version arXiv:2608.21855v1 (CC BY 4.0), distributed by arXiv under the Creative Commons Attribution 4.0 International licence, http://creativecommons.org/licenses/by/4.0/, as recorded in the arXiv licence metadata for this exact version and shown on the abstract page https://arxiv.org/abs/2608.21855v1. Full licence notice: \texttt{licenses/arxiv-2608.21855-CC-BY-4.0.txt}.\par\smallskip

\noindent\textit{Editorial scope.} Counted unit: the article's theorem thm-main (source0.tex lines 642--645). The article's complete original proof of it is delivered with it (source0.tex lines 646--685, one \texttt{proof} environment of 40 lines). No mathematics is written, repaired or re-derived: the statement and the proof are the article's own lines, in the article's own notation, and no retained line is substituted or re-flowed.  Retained uncounted as the source-local context the proof needs: lines 227--235; lines 604--611; lines 619--637.  Nothing was omitted from any retained span, so the package carries no omission record.  No retained line contains a citation, so the wrapper carries no bibliography and no bibliography entry.  No retained line contains a figure environment, an image-inclusion command or drawing code, so no image is delivered and the package declares no asset dependency.  Editorial formatting only, in addition: an AMS \texttt{amsart} preamble that restates the article's own declarations and macros used by the retained lines, namely the environments \texttt{definition}, \texttt{proposition}, \texttt{theorem} on separate counters, together with the standard packages those lines need.  The retained environments print in this document as Definition 1, Proposition 1, Theorem 1 in order of appearance, whereas the article prints the same items with the numbering of its own full document.  The complete native source of the article as distributed in the arXiv e-print for this exact version is bundled unchanged as \texttt{sources/208283/source0.tex}.
\begin{definition}
  \label{def-trapezoidal}
  Let $a_0,\dots,a_r$ be a strictly positive palindromic sequence. Put $m=\lfloor r/2\rfloor$ and $\delta_i=a_i-a_{i-1}$ for $1\leq i\leq m$.
  We say that the sequence is \emph{ trapezoidal} if
  \begin{equation}
    \delta_i\geq 0,\quad \delta_i=0 \Longrightarrow \delta_j=0\text{  for  } i\leq j\leq m.
  \end{equation}
  For a palindromic sequence, this is equivalent to the usual strictly increasing $\to$ constant $\to$ strictly decreasing formulation.
\end{definition}

\begin{proposition}
  \label{prop-bary-formula}
  For every circuit-generic $\rho$, one has the barycentric-fiber integration formula
  \begin{equation}
    f_{A,\rho}(t)=d! \int_{C_{\mathrm{generic}}} h(Q_x;t) d\mu(x).
  \end{equation}
  In particular, $f_{A,\rho}(t)$ is independent of $\rho$ and equals $f_A(t)$.
\end{proposition}

\begin{theorem}
  \label{thm-g-theorem}
  Let $Q$ be a simple $r$-polytope and
  \begin{equation}
    h(Q;t)=\sum_{i=0}^r h_i(Q)t^i.
  \end{equation}
  Then 
  \begin{equation*}
      h_0(Q)=1,\quad h_i(Q)=h_{r-i}(Q)\quad (0\leq i\leq r).
  \end{equation*}
  Moreover, the tuple $(g_0(Q),\dots,g_m(Q))$, $m=\lfloor r/2\rfloor$ defined by
  \begin{equation*}
    g_0(Q)=1,\quad g_i(Q)=h_i(Q)-h_{i-1}(Q), \,1\leq i\leq m
  \end{equation*}
  is an $M$-sequence. Consequently,
  \begin{equation}
    g_i(Q)\geq 0,\quad g_i(Q)=0 \Longrightarrow g_j(Q)=0\text{  for  } i\leq j\leq m.
  \end{equation}
\end{theorem}

\begin{theorem}
  \label{thm-main}
  Let $A$ be a full-dimensional flat vector arrangement. The coefficient sequence of $f_A(t)$ is trapezoidal.
\end{theorem}

\begin{proof}
  Let $r=N-d$.
  Since $\op{Ext}_\rho(B)\subset[N]\backslash B$ for every labeled basis $B$, one has $\deg f_A\leq r$.
  We can write
  \begin{equation*}
    f_A(t)=\sum_{i=0}^{r}a_i t^i,\quad h(Q_x;t)=\sum_{i=0}^r h_i(x)t^i,
  \end{equation*}
  where $h_i(x)$ is defined on $C_{\mathrm{generic}}$. Denote $m=\lfloor r/2\rfloor$.
  By Proposition \ref{prop-bary-formula}, we have
  \begin{equation*}
    a_i=d!\int_{C_{\mathrm{generic}}} h_i(x)d\mu(x).
  \end{equation*}
  
  By Theorem \ref{thm-g-theorem}, $h_i(x)=h_{r-i}(x)$, which gives
  \begin{equation}
    \label{eq-symmetry}
    a_i=a_{r-i},\quad i=0,\dots,m.
  \end{equation}
  
  For each $1\leq i\leq m$, 
  \begin{equation}
    \label{eq-monotone}
    a_i-a_{i-1}=d! \int_{C_{\mathrm{generic}}} (h_i(x)-h_{i-1}(x))d\mu(x)\geq0.
  \end{equation}
  The functions $h_i$ are measurable because, on $C_{\mathrm{generic}}$, they are coefficients of the finite sum $\tilde h_{A,\rho}$. Suppose $a_i-a_{i-1}=0$ for some $1\leq i\leq m$. Since the integrand in \eqref{eq-monotone} is nonnegative, $h_i(x)-h_{i-1}(x)=0$ for almost every $x\in C_{\mathrm{generic}}$.
  Therefore, Theorem \ref{thm-g-theorem} implies $h_j(x)-h_{j-1}(x)=0$ for almost every $x\in C_{\mathrm{generic}}$ for every $i\leq j\leq m$, and hence
  \begin{equation}
    \label{eq-plateau}
    a_j-a_{j-1}=0\text{  for  }i\leq j\leq m.
  \end{equation} 
  
  Finally, since $h_0(x)=1$ for $x\in C_{\mathrm{generic}}$,
  \begin{equation}
    \label{eq-positive}
    a_0=d! \mu(C_{\mathrm{generic}})=d! \mu(C)>0,
  \end{equation}
  where the last inequality holds because $C$ is full-dimensional in $H$. It follows from \eqref{eq-symmetry} and \eqref{eq-monotone} that $a_i>0$ for all $i$.
  
  All conditions of Definition \ref{def-trapezoidal} are now fulfilled. 
\end{proof}

\end{document}
